\section{Related Work}
\label{sec:related-work}

\paragraph{LLMs for cyber threat intelligence.}
LLMs have increasingly been applied to CTI workflows that require extracting, normalizing, and grounding security-relevant information from unstructured reports. A central task is mapping tactics, techniques, and procedures (TTPs) from natural-language threat reports to MITRE ATT\&CK. TRAM demonstrates an applied LLM pipeline for automated technique identification, motivated by the cost and brittleness of manual ATT\&CK mapping~\citep{ctid2023tram}. A recent systematization of automated TTP extraction methods, including generative LLMs, highlights persistent comparability challenges arising from heterogeneous ontologies, datasets, and evaluation protocols~\citep{buechel2025sok_ttp}. Expert-annotated resources such as AnnoCTR provide more controlled supervision and evaluation for ATT\&CK-labeled CTI text~\citep{lange2024annoctr}. Beyond report-level extraction, LLMs have been used to bootstrap structured CTI artifacts such as knowledge graphs~\citep{hu2024llm_tikg}, and hybrid knowledge-graph/LLM systems have been proposed for producing actionable intelligence from heterogeneous evidence~\citep{DBLP:conf/eurosp/FieblingerAR24}. Other work maps vulnerability descriptions into standardized taxonomies, where prompting alone remains unreliable but instruction templating and fine-tuning can improve grounding~\citep{liu2023not_end_of_story,zhang2024vtt_llm}. Instruction-tuned security models such as CyberPal and CyberPal~2.0 further improve cybersecurity-oriented behavior, while also illustrating the limitations of supervised fine-tuning for robust CTI reasoning and structured output generation~\citep{levi2024cyberpal,levi2025cyberpal2}. Complementary efforts study data-efficient ATT\&CK technique identification, including active learning~\citep{rahman2024alert}, and LLM pipelines for mapping detection artifacts, such as Sigma rules and SIEM analytics, to ATT\&CK via prompt chaining and retrieval~\citep{wudali2025ram}.

\paragraph{CTI and cybersecurity benchmarks.}
A growing set of benchmarks evaluates LLMs on CTI and cybersecurity workflows beyond general NLP tasks. CTIBench and AthenaBench target CTI-specific capabilities such as CVE$\rightarrow$CWE mapping, CVSS prediction, ATT\&CK technique extraction, mitigation recommendation, and threat-actor attribution~\citep{alam2024ctibench,alam2025athenabench}. SEvenLLM introduces a bilingual cybersecurity instruction corpus and benchmark with incident-analysis and response-oriented CTI tasks~\citep{ji2024sevenllm}. Broader cybersecurity evaluations include SECURE, which measures LLM performance across multiple cybersecurity tasks~\citep{DBLP:conf/acsac/BhusalANMLFFTBR24}; CyberMetric, which evaluates broad cybersecurity knowledge through 10{,}000 questions~\citep{cybermetric}; and CyberBench, which aggregates datasets for cybersecurity-language understanding tasks such as entity recognition, summarization, and classification~\citep{liu2024cyberbench}. These benchmarks expose recurring failure modes, including hallucination, mis-grounding, brittle identifier mapping, and schema errors. 

\paragraph{Reinforcement learning with verifiable rewards.}
Reinforcement learning with verifiable rewards (RLVR) has become a prominent post-training approach for tasks where correctness can be checked programmatically, especially mathematical reasoning and code generation~\citep{shao2024deepseekmath,deepseekr1}. Instead of relying on learned preference models, RLVR uses deterministic verifiers to assign rewards, making it attractive for domains with canonical answers and structured outputs. Recent work shows that RLVR can improve reasoning behavior, reliability, and calibration relative to supervised baselines~\citep{wen2025rlvr}, while other studies examine whether RLVR elicits new capabilities or primarily amplifies behaviors already present in the base model~\citep{yue2025does,cheng2025reasoning,wu2025invisibleleash}. Related analyses further suggest that supervised fine-tuning can overfit or memorize solution traces, whereas reinforcement learning may encourage more generalizable behavior under verifiable feedback~\citep{DBLP:conf/icml/ChuZYTXSLL025}. Follow-up work studies how RLVR gains depend on task difficulty, rollout budget, and exploration strategy~\citep{yang2025depth}, and shows that verifiable-reward formulations can transfer beyond math and code to other structured domains~\citep{lu2025generalization,su2025crossing}. Our work brings this paradigm to CTI, where many targets are canonical identifiers, label sets, or structured strings, and addresses the sparse-reward regime through hardness-gated answer-conditioned rationale generation and original-prompt distillation.
